\documentclass[10pt,a4paper]{article}
\usepackage[T1]{fontenc}
\usepackage[utf8]{inputenc}
\usepackage[ngerman]{babel}
\usepackage[a4paper,left=15.4mm,right=19.6mm,top=13.4mm,textheight=250mm]{geometry}
\usepackage{lmodern}
\usepackage{graphicx}
\usepackage[table]{xcolor}
\usepackage{tabularx,array,microtype,tikz}
\pagestyle{plain}
\setlength{\parindent}{0pt}
\setlength{\parskip}{0pt}
\renewcommand{\familydefault}{\sfdefault}
\definecolor{Navy}{HTML}{17324D}
\definecolor{Blue}{HTML}{3B6EA5}
\definecolor{Gold}{HTML}{F3B33D}
\definecolor{Ink}{HTML}{1E2935}
\definecolor{Muted}{HTML}{617080}
\definecolor{Line}{HTML}{888888}
\definecolor{Sky}{HTML}{EAF3FA}
\definecolor{Grid}{HTML}{BBBBBB}
\color{Ink}
\newif\ifsolution
\solutionfalse
% Gewählte Kurzarbeit: SA-Kopfanordnung, serifenlose Schrift und blaue Überschriften.
% Kopf nach Latex-SA-Vorlage/schulaufgabeVorlage.tex; keine fremden Rahmendaten.
\newcommand{\kopf}{\ifnum\value{page}=1
{\color{Blue}\fontsize{18}{21}\selectfont\bfseries\noindent\ifsolution Musterlösung\else Kurzarbeit\fi{} Informatik -- Datenbanken}\par
\vspace{1mm}\noindent Klasse 10\hspace{8mm}Datum: \rule{22mm}{.4pt}\hfill Name: \rule{45mm}{.4pt}\par\vspace{5mm}
\fi}
\newcommand{\sectiontitle}[1]{\vspace{1mm}{\color{Blue}\bfseries\large #1}\par\vspace{.3mm}{\color{Blue}\rule{\linewidth}{.7pt}}\par\vspace{1.5mm}}
\newcommand{\fuss}[1]{\par}
\newcommand{\gridfield}[2]{\par\vspace{1.5mm}\begin{tikzpicture}\path[use as bounding box] (0,0) rectangle (\linewidth,#1);\draw[Grid,line width=.25pt,step=5mm] (0,0) grid (\linewidth,#1);\draw[Line,line width=.45pt] (0,0) rectangle (\linewidth,#1);\ifsolution\node[anchor=north west,text width=\dimexpr\linewidth-6mm\relax,inner sep=0pt,font=\small,align=left,text=Blue] at (3mm,\dimexpr#1-2mm\relax) {#2};\fi\end{tikzpicture}\par}
\newcommand{\schematext}[1]{{\small\ttfamily #1}\par}
\newcommand{\dbtitle}[1]{{\small\bfseries\color{Blue}\texttt{#1}}\par\vspace{.7mm}}
\newcommand{\luecke}[1]{\ifsolution{\color{Blue}\bfseries #1}\else\rule{45mm}{.4pt}\fi}
\newcommand{\bewertung}[1]{\hfill{\normalsize[\hspace{8mm}/#1 BE]}}
% Zweigeteilte UML-Karte: Namensbereich (8 mm) und leerer Attributbereich (12 mm).
\newcommand{\klassenkarte}[6]{%
  \node[karte,minimum width=#4mm] (#1) at (#2,#3) {};
  \draw[Blue,line width=.7pt] ([yshift=2mm]#1.west) -- ([yshift=2mm]#1.east);
  \node[font=#6] at ([yshift=6mm]#1.center) {#5};}
\begin{document}
\kopf
\sectiontitle{1. Grundwissen:\bewertung{4}}
\textbf{Vervollständige} die Sätze, indem zu jedem Begriff aus dem Kontext Datenbanken den passenden Begriff aus dem Kontext Programmierung findest.\par\vspace{2mm}
Eine Spalte entspricht in der objektorientierten Programmierung \luecke{einem Attribut}.\par\vspace{3mm}
Eine Zeile entspricht in der objektorientierten Programmierung \luecke{einem Objekt}.\par\vspace{3mm}
Eine Tabelle entspricht in der objektorientierten Programmierung \luecke{einer Klasse}.\par\vspace{3mm}
Eine Zelle - au\ss{}er in der ersten Zeile - entspricht in der objektorientierten Programmierung\par\vspace{4mm}\luecke{einem Attributwert}.\par\vspace{3mm}
\sectiontitle{2. Apps auf dem Handy\bewertung{9}}
{\bfseries\color{Blue}Klassendiagramm Nr. 1}\par
\noindent\resizebox{\linewidth}{!}{\begin{tikzpicture}[x=1mm,y=1mm,karte/.style={draw=Blue,line width=.7pt,minimum height=20mm,outer sep=0pt}]
\path[use as bounding box] (0,0) rectangle (182,31);
\klassenkarte{app}{16}{16}{32}{App}{\bfseries}
\klassenkarte{handy}{91}{16}{32}{Handy}{\bfseries}
\klassenkarte{benutzer}{166}{16}{32}{Benutzer}{\bfseries}
\ifsolution
\draw[Blue,line width=.7pt] (app.east) -- node[pos=.12,above,font=\small] {n} node[pos=.88,above,font=\small] {m} node[below,font=\scriptsize] {ist installiert auf} (handy.west);
\draw[Blue,line width=.7pt] (handy.east) -- node[pos=.12,above,font=\small] {1} node[pos=.88,above,font=\small] {1} node[below,font=\scriptsize] {gehört zu} (benutzer.west);
\fi
\end{tikzpicture}}\par
\textbf{a) Ergänze} das Klassendiagramm, indem du Beziehungen (zwischen den Klassen App und Handy bzw. Handy und Benutzer) ergänzt. Wähle sinnvolle Kardinalitäten und benenne jede Beziehungen eindeutig.\par
\textit{Hinweis: Ein Benutzer hat in diesem Szenarion nur ein Handy und nicht mehrere.}\par\vspace{2mm}
\textbf{b) Benenne} die Beziehungen zwischen den Klassen und \textbf{begründe} deine Entscheidung.
\gridfield{35mm}{App--Handy: \textbf{n:m-Beziehung}. Eine App kann auf mehreren Handys installiert sein; ein Handy kann mehrere Apps enthalten.\par\smallskip Handy--Benutzer: \textbf{1:1-Beziehung}, wenn jedes Handy genau einem Benutzer gehört. Jeder Benutzer hat laut Szenario genau ein Handy. Die Gegenrichtung ist hier eine Modellannahme; andere passend begründete Modelle sind möglich.}
\vspace{2mm}
\textbf{c) Erkläre} bei welcher Beziehungsart eine Beziehungstabelle notwendig ist. \textbf{Begründe} deine Entscheidung.
\gridfield{30mm}{Eine Beziehungstabelle ist bei einer \textbf{n:m-Beziehung} notwendig. Ein einzelner Fremdschlüssel in einer der beiden Tabellen reicht für die vielen Zuordnungen in beiden Richtungen nicht aus. Jede Zeile der Beziehungstabelle verbindet über zwei Fremdschlüssel genau zwei Datensätze.}
\fuss{1}
\newpage
\kopf
\sectiontitle{3. Berechtigungen und Apps\bewertung{7}}
{\bfseries\color{Blue}Tabellenausschnitte Nr. 1}\par\vspace{1.5mm}
\begingroup\renewcommand{\arraystretch}{1.15}\setlength{\tabcolsep}{2mm}
\begin{minipage}[t]{.30\linewidth}\dbtitle{Berechtigung}
\begin{tabular}{|r|l|}\hline\rowcolor{Sky}\textbf{id}&\textbf{name}\\\hline
1&Kamera\\\hline 2&Standort\\\hline 3&Fotos lesen\\\hline
\end{tabular}
\end{minipage}\hfill
\begin{minipage}[t]{.27\linewidth}\dbtitle{Apps}
\begin{tabular}{|r|l|}\hline\rowcolor{Sky}\textbf{id}&\textbf{name}\\\hline
1&WhatsApp\\\hline 2&TikTok\\\hline 3&Instagram\\\hline
\end{tabular}
\end{minipage}\hfill
\begin{minipage}[t]{.39\linewidth}\dbtitle{Berechtigungen\_von\_Apps}
\begin{tabular}{|r|r|}\hline\rowcolor{Sky}\textbf{app\_id}&\textbf{berechtigung\_id}\\\hline
1&1\\\hline 1&3\\\hline 2&1\\\hline 3&2\\\hline 3&3\\\hline
\end{tabular}
\end{minipage}\par\endgroup\vspace{2mm}
{\bfseries\color{Blue}Klassendiagramm Nr. 2}\par
\noindent\resizebox{\linewidth}{!}{\begin{tikzpicture}[x=1mm,y=1mm,karte/.style={draw=Blue,line width=.7pt,minimum height=20mm,outer sep=0pt}]
\path[use as bounding box] (0,0) rectangle (182,46);
\klassenkarte{b}{18}{28}{36}{Berechtigungen}{\bfseries}
\klassenkarte{a}{168.5}{28}{27}{Apps}{\bfseries}
\draw[Blue,line width=.7pt] (b.east) -- node[pos=.04,above,font=\small] {n} node[pos=.96,above,font=\small] {m} node[below,font=\small] {ist freigegeben für} (a.west);
\end{tikzpicture}}\par
\textbf{a)} WhatsApp fragt beim Öffnen der App nach einer neuen Berechtigung \textquotedbl{}Standort\textquotedbl{}. Du nimmst die Berechtigungsanfrage an.\par
\textbf{Beschreibe} die Veränderungen, die passieren müssen, damit die Tabellen (siehe Tabellenausschnitte Nr. 1) auf dem aktuellen Stand sind.
\gridfield{30mm}{WhatsApp hat die \texttt{id} 1; Standort hat die \texttt{id} 2. In \texttt{Berechtigungen\_von\_Apps} wird eine neue Zeile mit \texttt{app\_id = 1} und \texttt{berechtigung\_id = 2} ergänzt. Beide Stammdatensätze existieren bereits; in \texttt{Apps} und \texttt{Berechtigung} sind keine neuen Zeilen erforderlich.}
\textbf{b) Formuliere} eine SQL-Anweisung, die zu jeder App alle passenden Berechtigungen, die freigegeben wurden, auflistet.
\gridfield{45mm}{{\ttfamily\fontsize{9}{11}\selectfont SELECT Apps.name, Berechtigung.name\\
FROM Apps, Berechtigung, Berechtigungen\_von\_Apps\\
WHERE Apps.id = Berechtigungen\_von\_Apps.app\_id\\
\hspace*{4mm}AND Berechtigung.id = Berechtigungen\_von\_Apps.berechtigung\_id;}\par\smallskip Jede freigegebene Zuordnung liefert eine Ergebniszeile mit Appname und Berechtigungsname. Nach 3a gehört auch WhatsApp--Standort dazu.}
\vspace{4mm}
\fuss{2}
\newpage
\kopf
\sectiontitle{4. Appentwickler\bewertung{5}}
Ein Appentwickler kann mehrere Apps entwickeln. Jede App wird in diesem Szenario von genau einem Appentwickler entwickelt.\par\vspace{2mm}
\begin{minipage}[t][18mm][t]{\linewidth}
\ifsolution
\schematext{App (id: int [PS], appentwickler\_id: int [FS],}
\schematext{\hspace{4mm}name: varchar(255), erscheinungsdatum: date,}
\schematext{\hspace{4mm}kategorie: varchar(255))}
{\small\color{Blue}\texttt{appentwickler\_id} verweist auf \texttt{Appentwickler.id}.}\par
\else
\schematext{App (id: int [PS], name: varchar(255),}
\schematext{\hspace{4mm}erscheinungsdatum: date, kategorie: varchar(255))}
\fi
\end{minipage}\par
{\small Verwendete Abkürzungen: PS = Primärschlüssel, FS = Fremdschlüssel.}\par\vspace{2mm}
\textbf{a) Beschreibe} folgende SQL-Anweisung in eigenen Worten.
\par\vspace{1mm}{\small\ttfamily\hspace{4mm}SELECT COUNT(*)\\
\hspace*{4mm}FROM Appentwickler\\
\hspace*{4mm}WHERE gruendungsjahr < 2025\par}
\gridfield{30mm}{Die Anweisung zählt die Appentwickler, die \textbf{vor dem Jahr 2025} gegründet wurden. Ausgegeben wird eine einzelne Zahl. Appentwickler mit Gründungsjahr 2025 werden nicht mitgezählt.}
\vspace{2mm}
\schematext{Appentwickler (id: int [PS], name: varchar(255), gruendungsjahr: int)}
\vspace{1mm}
\begin{tikzpicture}
\node[anchor=north west,inner sep=0pt,align=left,font=\small\ttfamily] (schema) at (0,0) {Zuordnung-Apps-Appentwickler (\\
\hspace*{4mm}app\_id: int [FS $\rightarrow$ App.id],\\
\hspace*{4mm}appentwickler\_id: int [FS $\rightarrow$ Appentwickler.id])};
\ifsolution
\draw[overlay,Blue,line width=.9pt] (schema.south west) -- (schema.north east);
\node[overlay,anchor=west,font=\small\bfseries,text=Blue] at ([xshift=4mm]schema.east) {entfällt};
\fi
\end{tikzpicture}\par
{\small In \texttt{Zuordnung-Apps-Appentwickler} bildet das Paar der beiden Fremdschlüssel den Primärschlüssel.}\par\vspace{2mm}
\textbf{b)} In den Tabellenschemata Apps, Appentwickler bzw. Zuordnung-Apps-Appentwickler haben sich Fehler eingeschlichen. \textbf{Korrigiere} den Fehler oder die Fehler. Du kannst direkt in das Tabellenschema hineinschreiben und Falsches durchstreichen, wenn notwendig.\par
\textbf{Begründe} deine Änderungen.
\gridfield{30mm}{Es liegt eine \textbf{1:n-Beziehung} zwischen Appentwickler und App vor. Die Zuordnungstabelle ist überflüssig und wird gestrichen. In \texttt{App} wird \texttt{appentwickler\_id: int [FS]} ergänzt; dieser Fremdschlüssel verweist auf \texttt{Appentwickler.id}. So ist jede App genau einem Entwickler zugeordnet. Das Schema \texttt{Appentwickler} bleibt unverändert.}
\vspace{2mm}
\textbf{c) Korrigiere} zwei Fehler, die in folgender SQL-Anweisung passiert sind. Schreibe die SQL-Anweisung darunter neu.
\par\vspace{1mm}{\fontsize{9.2}{11.2}\selectfont\ttfamily\hspace{4mm}SELECT\\
\hspace*{4mm}FROM Appentwickler, App\\
\hspace*{4mm}WHERE App.id = Appentwickler.id AND App.name = Appentwickler.name\par}
\gridfield{45mm}{{\ttfamily\fontsize{9}{11}\selectfont SELECT Appentwickler.name, App.name\\
FROM Appentwickler, App\\
WHERE App.appentwickler\_id = Appentwickler.id;}\par\smallskip Die Abfrage verwendet das in 4b korrigierte Schema. Eine mögliche Spaltenauswahl. Die fehlende SELECT-Liste wird ergänzt; die falsche Verknüpfung wird durch den Vergleich von Fremd- und Primärschlüssel ersetzt. Die zusätzliche Gleichheit der Namen entfällt.}
\par\vspace{2mm}{\bfseries Gesamt:\bewertung{25}}\par
\fuss{3}
\end{document}
